%% tox_dataset.tex -- shared sub-recipe. The step every other recipe assumes has
%% already happened: getting your data into TOX's arrays, and keeping it there.

\begin{recipe}[label=sub:tox-dataset]{Creating a TOX Dataset}

\recipepurpose{Read your expression table and gene families into the arrays
  Tensor Omics works on, check that they agree with each other, and store them
  together as a \file{.txdata} archive you can reopen later. Every other recipe
  starts from the arrays this one produces.}

\nomaintainer

\begin{requiredinputs}
  \item A gene $\times$ sample table (CSV/TSV): one identifier column, then one
        column per replicate or sample.
  \item A gene family table in OrthoFinder's \file{Orthogroups.tsv} layout: one
        family per row, one column per species, gene ids comma-separated.
  \item Nothing else --- this is the entry point.
\end{requiredinputs}

\begin{prerequisites}
  \item None. Later recipes depend on this one, not the other way round.
\end{prerequisites}

\begin{note}
  A TOX dataset is not an object. It is a handful of plain arrays that agree by
  \emph{position}: column $i$ of \code{expression} belongs to the gene named by
  \code{gene\_ids[i]}, and entry $i$ of \code{gene\_to\_family} names that
  gene's family in \code{family\_ids}. Nothing enforces this at run time, and
  nothing owns the arrays --- keeping them consistent is yours to do, which is
  what the validation step below is for.
\end{note}

\begin{workflow}
  \step{Measure your files first. The readers \emph{fill arrays you allocate},
        so the number of genes, the number of samples, the number of families
        and the width of the longest identifier all have to be known before the
        first read. Count them from the files.}
  \step{Read the identifiers with \code{read\_gene\_ids\_from\_tsv\_file},
        then the matrix with \code{read\_expression\_vectors\_tsv} into an
        \code{(n\_samples, n\_genes)} column-major array you allocated, then the
        families with \code{read\_orthofinder\_file}. A gene in no family gets
        the sentinel \code{0}.}
  \step{Check the arrays against each other before you rely on them:
        identifiers unique, family indices in range, expression non-negative.
        These are cheap and they catch the mistakes that are silent later.}
  \step{Store what you have with \code{save\_tox\_data}. Each array is passed
        together with the name it gets inside the archive; you may store any
        subset, so a dataset without centroids yet is perfectly valid.}
  \step{Reopen with \code{read\_tox\_data\_into}, which returns every member.
        Members that were never written come back with zero extents.}
\end{workflow}

\examplescript{python/examples/tox_dataset.py}{%
  Reads the bundled tables with the library's own readers, validates the
  relationships, writes an archive and reads it back to show the round trip is
  exact. Run unchanged from the repository root it takes about eight seconds
  over 88\,327 genes and writes \file{results/example.txdata}.}

\begin{codefragment}[language=Python,caption={Fragments to edit: the files, their shape, and what you store}]
import numpy as np
from tensor_omics import (read_gene_ids_from_tsv_file,
                          read_expression_vectors_tsv,
                          read_orthofinder_file, save_tox_data,
                          read_tox_data_into)

# --- sizes FIRST: the readers fill arrays you allocate -----
n_genes, n_samples, id_width  = 88327, 67, 14
n_families, family_id_width   = 15512, 9

gene_ids = read_gene_ids_from_tsv_file(
    "expression.tsv", id_width, n_genes,
    n_header_rows=1, gene_col=1)

# filled in place, so allocate it in its final form
expression = np.zeros((n_samples, n_genes),
                      dtype=np.float64, order="F")
read_expression_vectors_tsv(
    file_list=["expression.tsv"], gene_ids=gene_ids,
    expression_vectors=expression,
    n_header_rows=1, gene_col=1,
    value_cols=np.arange(2, n_samples + 2, dtype=np.int32),
    start_row=1)

families = read_orthofinder_file(
    "Orthogroups.tsv", gene_ids, family_id_width,
    n_families, n_genes)

# --- store: each array WITH the name it gets in the archive
save_tox_data("dataset.txdata",
    gene_ids=gene_ids,   gene_ids_file="gene_ids.bin",
    expression=expression, expression_file="expression.bin")

back = read_tox_data_into("dataset.txdata")
\end{codefragment}

\begin{parameters}
  \param{gene\_ids\_strlen}{Width of the fixed-size character slot each
    identifier is stored in}{Must be at least the longest identifier; measure
    it, do not guess. Anything longer is truncated, which silently breaks the
    match between the expression table and the family table}
  \param{value\_cols}{Which columns of the table hold values, 1-based}{%
    \code{2 \dots n\_samples+1} for the usual layout of one identifier column
    followed by the samples. Getting the count wrong leaves rows of the matrix
    untouched at zero}
  \param{n\_header\_rows}{Rows to skip before the data}{\code{1} for a single
    header line}
  \param{*\_file}{The name each array is given inside the archive}{Required
    whenever its array is passed, and vice versa. The names are yours to
    choose; the manifest records which is which}
\end{parameters}

\begin{expectedresults}
  On the bundled data the readers report 88\,327 genes $\times$ 67 samples and
  15\,512 families, of which 81\,960 genes are assigned to a family and 6\,367
  carry the unassigned sentinel. All four checks pass, and reading the archive
  back reproduces every array exactly. \code{get\_tox\_data\_dims} then reports
  \code{gene\_id\_len=14} and \code{family\_id\_len=9} --- the widths actually
  needed. The file is about 47\,MB for a 47.3\,MB matrix: the archive is a
  container, not compression, so expect roughly the size of the raw data.
\end{expectedresults}

\begin{pitfall}[Pitfall: sizing the arrays by guess]
  Every extent has to be supplied before the read, and a wrong one fails
  quietly rather than loudly. An identifier width below the longest id
  truncates identifiers, so genes stop matching between the two tables and the
  family assignment comes back mostly empty; a \code{value\_cols} shorter than
  the sample count leaves the remaining rows of the matrix at zero. Derive all
  of them from the files, as the example script does.
\end{pitfall}

\begin{pitfall}[Pitfall: passing an array without its member name]
  \code{save\_tox\_data} takes each array together with the name it gets inside
  the archive, and needs both. Passing only one half fails with
  \code{ToxInputError: invalid input arguments}, naming the missing argument.
  Do not work around it by dropping the array --- that is the data you meant to
  keep.
\end{pitfall}

\begin{pitfall}[Pitfall: assuming the archive validates itself]
  Nothing checks the arrays against each other on the way in or out. An
  archive whose \code{gene\_to\_family} is a different length from its
  \code{gene\_ids} stores and reloads without complaint, and the mismatch only
  surfaces as a wrong answer several recipes later. Validate before storing.
\end{pitfall}

\begin{pitfall}[Pitfall: reading a file that is not a TOX archive]
  The reader reports what went wrong at the file level, not at the schema
  level: a missing file gives \code{ToxIOError: could not open file}, an
  ordinary zip gives \code{Could not extract file from archive}, and an archive
  whose manifest lists nothing gives \code{could not read array data}. None of
  these means your analysis is wrong --- they mean the path is.
\end{pitfall}

\begin{interpretation}
  What you have at the end is the dataset every other recipe names in its
  required inputs. The expression matrix here is still \emph{raw}: normalize it
  before any distance or angle is computed from it, and store the normalized
  matrix rather than recomputing it, so that every later result refers to one
  known state of the data.

  Treat the archive as the unit of reproducibility. It records the arrays and
  their agreed positions, but not how they were produced --- not the
  normalization settings, not which centroid convention was used, not the
  outlier threshold. Keep that alongside it, because a \file{.txdata} on its own
  cannot tell you which of the choices in the recipes that follow were made.
\end{interpretation}

\begin{references}
  \reference{Emms \& Kelly (2019). OrthoFinder: phylogenetic orthology
    inference for comparative genomics. \emph{Genome Biology} 20:238.}
\end{references}

\end{recipe}
